

%\input{secao1}
\section{Fase Introdutória}
    No mês de abril e maio, o foco do projeto foi familiarização com ferramentas, nominalmente a estrutura básica de nós e cenas dentro da engine Godot (versão 4.5-dev3).
\subsection{Início}
    Foi atribuído aos membros do grupo a tarefa de buscar, assistir e seguir um tutorial de desenvolvimento de jogo na engine, onde não haveria cobrança de desenvolvimento de lógica profunda ou criação de biblioteca de recursos (imagens, cenários, etc.). 
    O repositório na qual os membros podem compartilhar seus avanços foi criado, e o projeto propriamente dito pode ter início. 
\subsection{Desenvolvimento}
    O plano, como segue no momento, é criar o mapa e iniciar movimentação básica na cena de exploração - ambiente na qual o personagem se desloca e pode encontrar inimigos. Dentre as tarefas: criar o mapa a partir de recursos obtidos online (TileMapLayer), criação do personagem e colisão (Area2D), implementação de movimentação (CharacterBody2D). Cabe ainda a atribuição de tarefas entre os membros. 

\section{Início do projeto}
    A partir do mês de junho, o projeto teve início propriamente dito. 
\subsection{Junho, Julho}
    Foi realizada a distribuição de tarefas para os membros do grupo. Entre elas, a criação da história, estabelecimento básico de mecânicas e sistemas de jogo, interação com NPCs e sistema de loja e inventário. Mesmo assim, pouco progresso propriamente dito foi feito.

\section{Inatividade/Downtime}
    Devido à falta de comunicação, diversos encargos acadêmicos, infortúnios técnicos e falta de organização, houve uma parada nos desenvolvimentos do projeto.
\subsection{Agosto, Setembro, Outubro}
    Algum progresso foi realizado. Foram desenvolvidos os menus de configurações, aprofundamento da história, mapa e movimentação, entre outros. Comunicação no grupo continuou precária. O objetivo para o próximo mês, e entrega final, é finalizar a cena de combate, menus de loja e inventário, o mapa final, e a integração dos módulos desenvolvidos por cada um dos membros.